\section{The published question and what its sources do not say}
\label{sec:open}

\subsection{The exact sentence and its referent}

The target sentence is \cite[Discussion, p.~i501]{shomorony2016}:

\begin{quote}
``The maximum-likelihood formulation of the AP (Medvedev and Brudno, 2009), on
the contrary, seems to be robust to these issues, and thus a good candidate for
the `correct' formulation. Understanding whether bridging conditions can be used
to guarantee that the maximum-likelihood sequence is the true sequence is
currently an open question.''
\end{quote}

A full-text scan of the accepted article finds no section, equation, formula, or
candidate-class pointer into Medvedev and Brudno. The sentence names the work
only by bibliography and the phrase ``maximum-likelihood formulation.'' The
surrounding paragraph contrasts information-feasible reconstruction with
``optimization-based formulation[s] of the AP such as those considered by
Nagarajan and Pop (2009) \dots\ and Medvedev and Brudno (2009)''
\cite{nagarajan2009,medvedev2009}; it fixes no objective and no tie rule. This is
the central source fact of this section: \srcfact{} the 2016 text does not
select a referent, a read-type convention, a candidate length, or a
maximizer/uniqueness reading.

\subsection{Four Medvedev--Brudno objects}

The cited paper contains four materially different objects, and they answer the
word ``sequence'' differently.

\begin{table}[t]
\centering
\small
\begin{tabular}{@{}p{0.30\textwidth}p{0.42\textwidth}p{0.20\textwidth}@{}}
\toprule
Reading & What the optimization returns & Sequence? \\
\midrule
(1) exact multinomial, candidate-intrinsic $N(D)$
  & a circular genome $D$
  & yes \\
(2) fixed-$N$ binomial approximation
  & a circular genome $D$ scored through per-type counts
  & yes \\
(3) Section~6.2 bidirected flow
  & a flow / ``(non-contiguous) assembly''
  & not necessarily \\
(4) broad ML principle
  & unspecified
  & n/a \\
\bottomrule
\end{tabular}
\caption{The four referents of ``the maximum-likelihood formulation
(Medvedev and Brudno, 2009).'' Reading (1) is what the paper \emph{names} as its
global read-count likelihood but explicitly abandons as not separable;
reading (2) is what its algorithm \emph{does} and what the contemporaneous
literature treats as the method; reading (3) is its algorithm's feasible set;
reading (4) is what the 2016 sentence most safely says literally. No primary
source selects one (\srcfact); the best operational sequence-level fit is (2)
(\srcinf).}
\label{tab:referents}
\end{table}

Reading (3) deserves emphasis: its feasible objects are flows, and Medvedev and
Brudno route the question of a single assembled sequence to a later heuristic.
Thus if the 2016 phrase meant the Section~6.2 flow literally, ``the
maximum-likelihood \emph{sequence}'' would be a category error. That is evidence
against reading (3) as the direct referent, while not excluding it as a
formalization of the underlying optimization.

Reading (2) is not a project invention. The dissertation version of the work
presents the fixed-size method as the operative one \cite{medvedev2010}, a
contemporaneous survey describes the Medvedev--Brudno assembler as requiring the
accurate genome size as a parameter \cite{howison2013}, and a follow-up
formulation advertises its improvement precisely as better genome-size handling
\cite{varma2011}. This is evidence about how the community read the method; it
does not upgrade reading (2) to the literal referent of the 2016 sentence.

A refinement matters for negative results. The binomial objective contains
$\log d_i$ and $\log(N-d_i)$, so it is defined only when $0<d_i<N$ for every
type; the exact multinomial is defined on every nonempty candidate. A refutation
therefore has to name the objective's domain, not just the candidate class.

\subsection{The axes the sentence leaves open}

\begin{table}[t]
\centering
\small
\begin{tabular}{@{}p{0.30\textwidth}p{0.36\textwidth}p{0.26\textwidth}@{}}
\toprule
Axis & Determination & Status \\
\midrule
named likelihood layer & readings (1)--(4) all literal & \srcgap \\
read-type space & oriented windows vs.\ molecule classes & \srcgap \\
candidate length & not fixed by the sentence; fixed length is best
  operational fit & \srcgap \\
maximizer vs.\ uniqueness & both literal & \srcgap \\
cyclic shift in $\approx$ & \textbf{forced} for the strong schema & \mathfact \\
reverse complement in $\approx$ & \textbf{forced iff} molecule read types
  & \mathfact\ + \srcgap \\
other genome equivalences & no source support & \srcfact \\
tie rule & none stated; many optima possible & \srcfact \\
\bottomrule
\end{tabular}
\caption{What the published sentence does and does not determine. The two
``forced'' rows are mathematical consequences rather than readings; the
remaining axes are genuine source gaps.}
\label{tab:axes}
\end{table}

Three points from Table~\ref{tab:axes} are worth isolating.

\paragraph{Candidate length is not the same as known genome size.}
Medvedev and Brudno's Section~6.1 approximation supplies an external genome-size
parameter $N$ to the likelihood. It does \emph{not} thereby restrict every
candidate $D$ to $|D|=N$. Knowing or supplying $N=|S|$ and restricting
candidates to $|D|=|S|$ are different assumptions; a same-length theorem is a
restricted result unless a source argument fixes the restriction.

\paragraph{Cyclic shift is mandatory, reverse complement is coupled.}
By Remark~\ref{rem:cyclic} the strong schema requires cyclic shift in $\approx$.
Reverse complement is forced precisely when read types are molecule classes: a
molecule and its reverse complement are always tied under class-indexed
observations, while under oriented reads the two count vectors generally differ.
Hence ``oriented with cyclic shift'' and ``molecular with dihedral equivalence''
are the two coherent source-supported panels; they are not freely mixable.

\paragraph{Ties are generic.}
The exact circular likelihood is invariant under cyclic shift, so distinct
spellings tie. A contemporaneous formulation states that many Eulerian tours can
have equal likelihood, so ``the maximum-likelihood sequence'' cannot be assumed
unique without an equivalence convention.

\subsection{Residual source gaps}

We record two gaps that bear on how strongly any formal result can be attributed
to the published sentence. First, the accepted publisher supplement of
\cite{shomorony2016} --- the last unexamined accepted artifact that could contain
a likelihood definition --- has not been retrieved through the available
execution path; a definition located there could narrow the referent. Second,
Medvedev and Brudno's Section~6.1 writes that there are ``$4^k$ such variables''
while its model vocabulary is $k$-molecules (reverse-complement classes), so the
index set and the vocabulary are in internal tension. A formalization must name
which it uses. Neither gap changes the mathematics below; both change how
directly a result may be advertised as settling the historical question.

\paragraph{Consequence for this paper.}
Because no referent is selected, we do not present any single theorem as ``the''
answer to the 2016 question. Instead, Section~\ref{sec:finite} proves and refutes
statements for named points in the axis space, and Section~\ref{sec:population}
presents a project-level repaired model that is explicitly not a historical
claim. A formal statement counts as a settlement of the published problem only
after an argument fixes the axes of Table~\ref{tab:axes} to a source-supported
panel; that argument is not currently available.
